\FloatBarrier
\newpage
\section{积分的计算}
\label{sec:calculus-integration-methods}

\begingroup
\setlength{\parskip}{3pt}
\setlength{\abovedisplayskip}{10pt}
\setlength{\belowdisplayskip}{10pt}

基本定理把定积分的计算与原函数联系起来。本节由求导结构组织寻找原函数的方法，再回到分割求和，研究没有现成原函数表达式时的数值近似。

\subsection{不定积分}

\begin{definition}
函数\(f\)在一个区间上的全部原函数，称为它在该区间上的{\heiti{}不定积分}，记为\(\int f(x)\,dx\)。
\end{definition}

若\(F'=f\)，则\(F+C\)也是原函数。反过来，任意两个原函数之差的导数为零，由中值定理，它们只相差常数。因此
\begin{equation*}
 \int f(x)\,dx=F(x)+C,
\end{equation*}
其中\(C\)为任意常数。定积分是确定的数，不定积分表示一个函数族；写出其中一个原函数，就能通过加常数得到全部原函数。

将基本求导公式反过来使用，得到
\begin{equation*}
 \begin{aligned}
 \int x^\alpha\,dx&=\frac{x^{\alpha+1}}{\alpha+1}+C\quad(\alpha\ne-1),&
 \int\frac{dx}{x}&=\ln|x|+C,\\
 \int e^x\,dx&=e^x+C,&\int a^x\,dx&=\frac{a^x}{\ln a}+C,\\
 \int\sin x\,dx&=-\cos x+C,&\int\cos x\,dx&=\sin x+C,\\
 \int\frac{dx}{1+x^2}&=\arctan x+C,&
 \int\frac{dx}{\sqrt{1-x^2}}&=\arcsin x+C,\\
 \int\frac{dx}{\cos^2x}&=\tan x+C,&
 \int\frac{dx}{\sin^2x}&=-\cot x+C,\\
 \int\sinh x\,dx&=\cosh x+C,&\int\cosh x\,dx&=\sinh x+C.
 \end{aligned}
\end{equation*}
各式沿用相应函数的定义范围，\(a>0\)、\(a\ne1\)。涉及\(1/x\)时，在不跨过零点的区间上使用；两个不同区间上的积分常数可以独立选择。

由求导的线性性质，不定积分也可以按线性组合寻找原函数。核对结果最直接的方法，是对所得表达式求导。

\begin{example}
求\(\int(3x^2-2\cos x+1/x)\,dx\)。再求满足\(F'(x)=2x+e^x\)、\(F(0)=3\)的函数\(F\)，说明初值怎样确定常数。
若另要求\(F(1)=5\)，这些条件能否同时满足？
再求\(\int(x+1)^2\,dx\)与\(\int(1+x^2)^{-1}\,dx\)，比较先作代数展开与直接识别基本公式的选择。
\end{example}

\subsection{换元积分}

若\(F'=f\)，链式法则给出
\begin{equation*}
 \frac d{dx}F(u(x))=f(u(x))u'(x).
\end{equation*}
因此，将内层函数的导数与外层表达式一同识别，就得到换元积分：
\begin{equation*}
 \int f(u(x))u'(x)\,dx=F(u(x))+C.
\end{equation*}
微分记号\(du=u'(x)\,dx\)将它写为\(\int f(u)\,du\)，依据是链式法则及微分形式不变性。

例如，在\(\int2x\cos(x^2)\,dx\)中令\(u=x^2\)，则\(du=2x\,dx\)，得到\(\sin(x^2)+C\)。若缺少这个\(2x\)，就不能把同一替换直接当作完整的换元。

\begin{theorem}[定积分换元公式]\label{thm:calculus-substitution}
设\(f\)在包含\(\varphi([\alpha,\beta])\)的区间上连续，\(\varphi\)在\([\alpha,\beta]\)上连续可导，则
\begin{equation}\label{eq:calculus-substitution}
 \int_\alpha^\beta f(\varphi(t))\varphi'(t)\,dt
 =\int_{\varphi(\alpha)}^{\varphi(\beta)}f(x)\,dx.
\end{equation}
\end{theorem}

\begin{proof}
取\(f\)的原函数\(F\)，左端由链式法则及Newton--Leibniz公式等于\(F(\varphi(\beta))-F(\varphi(\alpha))\)，也就是右端。
\end{proof}

这个公式不要求\(\varphi\)单调。例如\(\varphi(t)=\sin t\)、\(0\le t\le\pi\)时，代换后的起点和终点相同，有向贡献在往返中相消。若把一个给定积分改用新参数计算，则应选好对应的参数范围，并同时变换微分和上下限。

\begin{example}
计算\(\int_0^1 2x e^{x^2}\,dx\)、\(\int_0^1 dx/(1+x)^2\)及\(\int_{-\pi/2}^{\pi/2}\cos^2t\,dt\)。再用\(x=R\sin t\)计算\(\int_{-R}^R\sqrt{R^2-x^2}\,dx\)，与半圆面积比较。
\end{example}

\begin{exercise}
\mbox{}\par\nobreak

（1）求\(\int x/(1+x^2)\,dx\)、\(\int e^{2x+1}\,dx\)和\(\int\cos x/(1+\sin^2x)\,dx\)。

（2）计算\(\int_0^\pi f(\sin t)\cos t\,dt\)，其中\(f\)在\([0,1]\)上连续。解释结果与参数往返的关系。

（3）在\(\int_1^e(\ln x)^2/x\,dx\)中令\(u=\ln x\)，写出新积分的上下限并计算。

（4）求\(\int(\sqrt{x}+1/\sqrt{x})\,dx\)，其中\(x>0\)。再求满足\(F'(x)=3x^2-2x+1\)、\(F(1)=0\)的原函数，说明两问中积分常数的作用。

（5）计算\(\int x\sqrt{1+x^2}\,dx\)与\(\int_0^1x/(1+x^2)^2\,dx\)，指出所选中间变量及对应微分。

（6）计算\(\int_1^e dx/[x(1+\ln^2x)]\)与\(\int_0^{\pi/4}\tan x\,dx\)，说明所选的中间变量。
\end{exercise}

\subsection{分部积分}

乘积法则\((uv)'=u'v+uv'\)给出
\begin{equation}\label{eq:calculus-integration-by-parts}
 \int_a^b u(x)v'(x)\,dx
 =\left.u(x)v(x)\right|_a^b-\int_a^b u'(x)v(x)\,dx,
\end{equation}
其中\(u,v\)连续可导。不定积分的相应形式为\(\int u\,dv=uv-\int v\,du\)。

选择分部的目的，是让新的积分更容易处理。例如，对\(\int x e^x\,dx\)，将\(x\)求导、将\(e^x\)积分，得到
\begin{equation*}
 \int x e^x\,dx=xe^x-\int e^x\,dx=(x-1)e^x+C.
\end{equation*}
多项式因子求导后次数降低，因而这一选择可以继续使用。对\(\ln x\)，则可以把它看成\(\ln x\cdot1\)，通过积分\(1\)引入\(x\)。

对\(m\)次多项式\(P\)，连续分部直到多项式因子成为常数，得到
\begin{equation*}
 \int P(x)e^x\,dx
 =e^x\bigl(P(x)-P'(x)+\cdots+(-1)^mP^{(m)}(x)\bigr)+C.
\end{equation*}
对右端求导时，括号内的每一项都分别产生自身与下一阶导数，相邻项恰好抵消，只留下\(P(x)e^x\)。多项式求导最终为零，使分部过程有限终止；这种次数的降低，比只看表达式含有哪些因子更能说明选择的依据。

分部后也可能重新出现原积分。这时可以先计算一个定积分并记为\(J\)，将重复出现的\(J\)移到同一边，通过一个关于\(J\)的方程求值。对不定积分也可作相应整理，最后用求导核对，并保留任意常数。

\begin{example}
求\(\int\ln x\,dx\)、\(\int x\sin x\,dx\)及\(\int_0^1 x^2e^x\,dx\)。再计算\(\int e^x\cos x\,dx\)，观察两次分部后原积分如何重新出现。
\end{example}

\begin{exercise}
\mbox{}\par\nobreak

（1）计算\(\int_0^1x\arctan x\,dx\)及\(\int_1^e x\ln x\,dx\)。

（2）记\(J_n=\int_0^1x^ne^x\,dx\)。建立\(J_n\)与\(J_{n-1}\)的关系，并求\(J_3\)。

（3）设\(F\)在\([0,1]\)上连续可导，记\(f=F'\)。证明\(\int_0^1F(x)\,dx=F(1)-\int_0^1xf(x)\,dx\)。

（4）计算\(\int_0^1x\ln(1+x)\,dx\)。分部后得到的有理式可以先做怎样的代数整理？

（5）记\(I_n=\int_0^{\pi/2}\sin^nx\,dx\)。对\(n\ge2\)，用分部积分建立\(I_n\)与\(I_{n-2}\)的关系，并求\(I_4\)。

（6）分别计算\(\int xe^{x^2}\,dx\)与\(\int xe^x\,dx\)，比较换元与分部的选择。两个表达式都含有因子\(x\)，为什么不宜据此选择同一种方法？
\end{exercise}

\newpage
\subsection{定积分的计算}

定积分有时可以在寻找原函数之前，先由区间与函数的关系简化。作\(x=-t\)的换元，得到
\begin{equation*}
 \int_{-a}^a f(x)\,dx=\int_0^a(f(x)+f(-x))\,dx.
\end{equation*}
所以奇函数在对称区间上的积分为零，偶函数的积分为正半段积分的两倍。

同样，在\([a,b]\)上作\(x=a+b-t\)，有
\begin{equation}\label{eq:calculus-integral-reflection}
 \int_a^b f(x)\,dx=\int_a^b f(a+b-x)\,dx.
\end{equation}
若\(f(x)+f(a+b-x)=c\)，将两式相加就得到\(\int_a^b f=c(b-a)/2\)，无需求出原函数。

例如，\(f(x)=1/(1+e^x)\)满足\(f(x)+f(-x)=1\)，故\(\int_{-a}^a dx/(1+e^x)=a\)。这个计算使用了函数在整个区间上的配对关系。

若\(f(a+b-x)=f(x)\)，记\(J=\int_a^bxf(x)\,dx\)，反射换元给出\(J=\int_a^b(a+b-x)f(x)\,dx\)。相加即得
\begin{equation*}
 \int_a^bxf(x)\,dx=\frac{a+b}{2}\int_a^bf(x)\,dx.
\end{equation*}
对称关系使一次因子\(x\)可以由区间中点代替。这里同样没有要求先求出\(xf(x)\)的原函数。

绝对值、分段表达式和几何面积则通常需要先拆分区间，确定每段上的表达式或符号，再使用基本定理。

\begin{example}
计算\(\int_{-2}^2(x^3+\cos x)\,dx\)、\(\int_0^2|x^2-1|\,dx\)及\(\int_0^{\pi/2}\sin x/(\sin x+\cos x)\,dx\)。分别指出计算中起主要作用的关系。
\end{example}

\begin{exercise}
\mbox{}\par\nobreak

（1）设\(f\)连续且周期为\(T>0\)。证明\(\int_c^{c+T}f\)与\(c\)无关。可以拆分区间，也可以把它看成变限积分求导。

（2）设\(f\)在\([0,1]\)上可积，且\(f(x)+f(1-x)=x^2+(1-x)^2\)，求\(\int_0^1f(x)\,dx\)。

（3）计算\(\int_{-1}^2|x(x-1)|\,dx\)，并区分结果与\(\int_{-1}^2x(x-1)\,dx\)。

（4）用\(x=\pi-t\)的换元计算\(\int_0^\pi x\sin x\,dx\)，再用分部积分核对。比较两种方法分别利用了什么结构。

（5）不先求原函数，计算\(\int_0^1\dfrac{\ln(1+x)}{\ln(1+x)+\ln(2-x)}\,dx\)。说明\(x\)与\(1-x\)处的函数值怎样配对。
\end{exercise}

\subsection{数值积分}

即使不能从基本公式中直接找到原函数，也可以回到Riemann和近似积分。对\([a,b]\)作\(n\)等分，令\(h=(b-a)/n\)、\(x_i=a+ih\)，常用的近似有
\begin{equation*}
 L_n=h\sum_{i=0}^{n-1}f(x_i),\qquad
 R_n=h\sum_{i=1}^nf(x_i),\qquad
 M_n=h\sum_{i=1}^nf\!\left(\frac{x_{i-1}+x_i}{2}\right).
\end{equation*}
它们依次为左、右端点与中点矩形公式。各段改用弦近似，得到{\heiti{}梯形公式}
\begin{equation*}
 T_n=h\left(\frac{f(a)+f(b)}2+\sum_{i=1}^{n-1}f(x_i)\right)
 =\frac{L_n+R_n}{2}.
\end{equation*}
连续性保证这些近似都趋于积分，但固定\(n\)时的误差还取决于函数的变化。

若\(f\)单调递增，则\(L_n\le\int_a^b f\le R_n\)，且
\begin{equation*}
 R_n-L_n=h(f(b)-f(a)).
\end{equation*}
单调递减时上下界交换。若\(f\)为凸函数，图像不高于弦，故积分不超过梯形和；中点切线不高于图像，且其斜率项在对称小段上积分为零，故
\begin{equation*}
 M_n\le\int_a^b f(x)\,dx\le T_n.
\end{equation*}
凹函数的方向相反。

\begin{figure}[!ht]
\centering
\includegraphics[width=.78\linewidth]{calculus-quadrature}
\caption{凸函数的中点与梯形估计。左图各段中点切线的积分等于对应矩形面积，右图用弦给出上界。}
\label{fig:calculus-quadrature}
\end{figure}

若函数的单调性或凹凸性在区间内改变，可以先在变化点拆分，再对各段分别建立上、下界，最后相加。所需信息由当前的估计方法决定：单调性提供矩形界，凹凸性提供中点与梯形界；仅有有限个取样值时，不能自动获得这些界的方向。

若\(A\le I\le B\)，以\((A+B)/2\)近似\(I\)，误差至多为\((B-A)/2\)。故两界之差不超过\(2\varepsilon\)时，该近似的误差不超过\(\varepsilon\)。

\FloatBarrier
\newpage

将单段的中点与梯形估计用于\(e^t\)在\([\ln a,\ln b]\)上的积分，就得到
\begin{equation*}
 \sqrt{ab}\le\frac1{\ln b-\ln a}\int_{\ln a}^{\ln b}e^t\,dt
 =\frac{b-a}{\ln b-\ln a}\le\frac{a+b}{2}.
\end{equation*}
因指数函数严格凸，两侧均取严格不等号。这重新解释了命题\ref{prop:calculus-logarithmic-mean}中的平均值关系：对数平均也是指数函数在相应区间上的积分平均值。前面用导数建立的比较，在累积量中获得了同一个来源。

\begin{example}
对\(J=\int_0^1e^{-x^2}\,dx\)，写出四等分的左右端点估计，给出一个包含\(J\)的区间。再对\(\int_1^2dx/x\)写出两等分的中点与梯形估计，利用凸性判断它们的方向。
\end{example}

\begin{exercise}
\mbox{}\par\nobreak

（1）对\(\int_0^1x^2\,dx\)，求\(M_n,T_n\)及其精确误差，比较误差随\(n\)的变化。

（2）只使用\(e^{-x^2}\)的单调性，要使左右矩形和之差小于\(10^{-2}\)，可以取怎样的分段数？

（3）一个连续函数的端点值相同，是否足以断定梯形公式精确？举例说明。

（4）对\(\int_0^1\ln(1+x)\,dx\)，写出两等分的中点与梯形近似，利用凹性判断上下界，再计算积分值核对。

（5）已知\(f\)在\([0,1]\)上连续、单调递减，且\(f(0)=2\)、\(f(1/2)=1\)、\(f(1)=1/2\)。仅由这些数据给出积分的上下界，再写出梯形近似。这个近似是否必为精确值？

（6）设\(f(x)=x^3\)，\(-1\le x\le2\)。在零点拆分区间，各用一个中点矩形和一个梯形，分别求整个积分的上下界。若直接在\([-1,2]\)上只用一个梯形，它与真实积分的大小能否仅由整段凸性判断？

（7）用两等分的中点与梯形公式估计\(\int_0^1\sqrt{1+x}\,dx\)。取两界平均值作近似，给出无需计算精确积分的误差界。

（8）用梯形公式近似\(\ln2=\int_0^1dx/(1+x)\)。只利用单调性给出误差界，并确定使误差不超过\(10^{-3}\)的一个充分分段数。若只取左端点和，所需的分段数有何变化？

（9）设\(f\)在\([0,2]\)上连续且为凸函数，已知\(f(0)=1\)、\(f(1)=2\)、\(f(2)=5\)。用一个中点矩形和一个梯形给出积分的上下界，再改用两个梯形求上界。若还知\(f(1/2)=5/4\)、\(f(3/2)=13/4\)，写出两个中点矩形给出的下界，并估计两界平均值的误差。

\end{exercise}

\endgroup
\FloatBarrier
\newpage
\begingroup
\setlength{\parskip}{4pt}
\setlength{\abovedisplayskip}{10pt}
\setlength{\belowdisplayskip}{10pt}
\subsection{Taylor公式的积分余项*}

分部积分还可以把局部展开的余项写成累积量。先看一次近似：由基本定理，\(f(x)-f(a)=\int_a^xf'(t)\,dt\)。若进一步把\(f'\)的变化纳入积分，就能把线性项分离出来；重复这一过程，便得到任意有限阶的积分余项。

\begin{theorem}\label{thm:calculus-taylor-integral}
设\(f\)在包含\(a,x\)的区间上具有直到\(n+1\)阶的连续导数，则
\begin{equation}\label{eq:calculus-taylor-integral-remainder}
 f(x)=\sum_{k=0}^n\frac{f^{(k)}(a)}{k!}(x-a)^k
 +\frac1{n!}\int_a^x(x-t)^nf^{(n+1)}(t)\,dt.
\end{equation}
\end{theorem}

\begin{proof}
记右端积分项为\(Q_n\)。对\(n\ge1\)，分部积分得到
\begin{equation*}
 Q_n=-\frac{f^{(n)}(a)}{n!}(x-a)^n
 +\frac1{(n-1)!}\int_a^x(x-t)^{n-1}f^{(n)}(t)\,dt
 =Q_{n-1}-\frac{f^{(n)}(a)}{n!}(x-a)^n.
\end{equation*}
而\(Q_0=\int_a^xf'(t)\,dt=f(x)-f(a)\)。逐次代回，即得公式。
\end{proof}

记\(h=x-a\)，在余项中作\(t=a+sh\)的换元，可将不同长度的积分区间统一到\([0,1]\)：
\begin{equation}\label{eq:calculus-taylor-scaled-remainder}
 Q_n=\frac{h^{n+1}}{n!}\int_0^1(1-s)^n f^{(n+1)}(a+sh)\,ds.
\end{equation}
这个表示对正、负\(h\)都成立。因子\(h^{n+1}\)记录误差的阶数，其余部分是高阶导数沿基点与终点之间的加权累积。

若\(|f^{(n+1)}|\le M\)，由积分绝对值估计得到
\begin{equation*}
 |Q_n|\le\frac{M}{(n+1)!}|x-a|^{n+1}.
\end{equation*}
因为\(\int_0^1(1-s)^n\,ds=1/(n+1)\)，上界中的阶乘正是对这个权函数积分后得到的。若高阶导数在不同位置大小差别明显，保留积分表达式还可以得到比统一上界更细的估计。

\paragraph{与Lagrange余项的联系} \mbox{}

权函数\((1-s)^n\)非负。设\(f^{(n+1)}\)在所考察线段上的最小值、最大值为\(m,M\)，则
\begin{equation*}
 \frac m{n+1}\le\int_0^1(1-s)^nf^{(n+1)}(a+sh)\,ds\le\frac M{n+1}.
\end{equation*}
由连续性与介值定理，可在\(a,x\)之间取\(\xi\)，使中间积分等于\(f^{(n+1)}(\xi)/(n+1)\)。代回式\eqref{eq:calculus-taylor-scaled-remainder}，便得到
\begin{equation*}
 Q_n=\frac{f^{(n+1)}(\xi)}{(n+1)!}(x-a)^{n+1}.
\end{equation*}
Lagrange余项将这一段高阶变化的加权平均，由某一点的值实现；积分余项则保留各位置对误差的贡献。

\newpage
\paragraph{局部阶数与余项符号} \mbox{}

当\(n\ge1\)时，若\(f\)具有连续的\(n\)阶导数，对\(n-1\)阶积分展开再减去第\(n\)项，可得
\begin{equation*}
 f(a+h)-T_n(a+h)=\frac{h^n}{(n-1)!}
 \int_0^1(1-s)^{n-1}\bigl(f^{(n)}(a+sh)-f^{(n)}(a)\bigr)\,ds.
\end{equation*}
由\(f^{(n)}\)在\(a\)处连续，括号内的差在整段\(0\le s\le1\)上一致趋零，故右端为\(o(h^n)\)。这在当前连续高阶导数条件下再次得到Peano形式；5.5中的Peano定理还允许更弱的条件。

对一次近似，式\eqref{eq:calculus-taylor-scaled-remainder}给出特别直观的关系：
\begin{equation*}
 f(a+h)-f(a)-f'(a)h=h^2\int_0^1(1-s)f''(a+sh)\,ds.
\end{equation*}
若\(m\le f''\le M\)，则无论\(h\)的正负，都有
\begin{equation}\label{eq:calculus-quadratic-tangent-bounds}
 \frac m2h^2\le f(a+h)-f(a)-f'(a)h\le\frac M2h^2.
\end{equation}
二阶导数非负时，图像位于切线上方；二阶导数还有正下界时，图像与切线的距离至少含一个确定的二次贡献。这也接回了5.5的凸性与习题中的定量估计。

\begin{example}
设\(f(x)=\ln(1+x)\)，\(0\le x\le1\)。写出一次近似的积分余项，并由\(f''\)的上下界估计\(f(x)-x\)。再对\(e^x\)使用同一关系，比较这两个函数与零点切线的位置。
\end{example}

\begin{exercise}
\mbox{}\par\nobreak

（1）用积分余项证明\(x>0\)时\(e^x>1+x+x^2/2\)，与Lagrange余项的证明比较。

（2）对\(f(x)=\ln(1+x)\)，写出在零点的二阶积分余项，证明\(x\ge0\)时\(x-x^2/2\le\ln(1+x)\le x-x^2/2+x^3/3\)。

（3）设\(f''\ge m>0\)。对区间中点的左右增量分别使用式\eqref{eq:calculus-quadratic-tangent-bounds}，重新证明习题中\(\bigl(f(a)+f(b)\bigr)/2-f((a+b)/2)\)的下界。

（4）设\(f\)在\([0,1]\)上三阶连续可导。用积分余项比较\(f(1)\)与二阶Taylor多项式在\(1\)处的值。当\(f'''\)在前半段与后半段大小不同时，权函数如何影响两部分的贡献？

（5）设\(f'''\ge0\)且连续。用积分余项判断\(f(x)\)与\(T_2(x)\)在基点\(a\)左右两侧的大小关系，解释为什么两侧方向不同。

\begin{samepage}
（6）设\(f(0)=1\)、\(f'(0)=0\)，且在\([0,1]\)上\(f''\)连续并满足\(2\le f''\le4\)。先估计\(f(x)\)，再给出\(\int_0^1f(x)\,dx\)的上下界。
\par
\end{samepage}
\end{exercise}

\subsection{数值积分的误差估计*}

中点公式的线性误差可以在每个小段内相消，因此比单纯估计函数波动更精确。下面只研究这一种公式。

\begin{proposition}\label{prop:calculus-midpoint-error}
若\(f\)在\([a,b]\)上二阶连续可导，且\(|f''|\le M\)，则\(n\)等分的中点公式满足
\begin{equation*}
 \left|\int_a^b f(x)\,dx-M_n\right|\le\frac{M(b-a)^3}{24n^2}.
\end{equation*}
\end{proposition}

\begin{proof}
对长度为\(h\)、中点为\(m\)的小段，Taylor公式给出
\begin{equation*}
 |f(x)-f(m)-f'(m)(x-m)|\le\frac M2(x-m)^2.
\end{equation*}
线性项在\([m-h/2,m+h/2]\)上的积分为零。因此该小段的误差不超过
\begin{equation*}
 \frac M2\int_{m-h/2}^{m+h/2}(x-m)^2\,dx=\frac{Mh^3}{24}.
\end{equation*}
把\(n\)段误差相加，再代入\(h=(b-a)/n\)，便得结论。
\end{proof}

证明先用局部多项式近似，再把各段误差累积。真正进入最终上界的是整个区间上统一有效的\(M\)，而每一段线性项的相消解释了\(n^{-2}\)的误差阶。

若分割不等长，同一证明给出\(\frac M{24}\sum_i(\Delta x_i)^3\)的误差界；若各段使用各自的二阶导数上界\(M_i\)，还可以写为\(\frac1{24}\sum_iM_i(\Delta x_i)^3\)。因此函数弯曲较明显的位置可以分得更细，平缓部分则不必采用同样的段长。

又因为每段中点切线的线性项积分为零，当\(f''\ge0\)时，中点近似不超过真实积分。这将数值近似的方向、误差大小与前面凸性的几何判断联系起来。

\begin{exercise}
\mbox{}\par\nobreak

（1）对\(f(x)=e^{-x^2}\)、\(0\le x\le1\)，验证\(|f''(x)|\le2\)。要使中点公式误差小于\(10^{-4}\)，可以选取怎样的\(n\)？

（2）将上述证明用于\(f(x)=x^2\)，比较误差界与精确误差。为什么在这个例子中等号能够成立？

\begin{samepage}
（3）对\(\int_0^1x^2\,dx\)，分别使用分割\(0,1/2,3/4,1\)与三等分分割的中点近似，求精确误差，并与\(\frac M{24}\sum_i(\Delta x_i)^3\)核对。相同分段数下，哪一个分割更准确？
\par
\end{samepage}

（4）若\(f''\)连续且\(0<m\le f''\le M\)，给出任意分割的中点近似误差的正下界与上界。说明二阶导数的符号和大小分别提供什么信息。
\end{exercise}

\endgroup
\newpage
\subsection{特殊换元*}

二次表达式可以通过配方、因式分解或三角参数转为已知形式。方法的选择取决于当前表达式具有的结构。

对\(a>0\)，在\(\sqrt{a^2-x^2}\)中令\(x=a\sin t\)、\(-\pi/2\le t\le\pi/2\)，根式化为\(a\cos t\)。相应的\(dx=a\cos t\,dt\)仍须同时代入。类似地，\(x=a\tan t\)可以处理\(a^2+x^2\)，也可由双曲函数恒等式选择\(x=a\sinh t\)。

有理式则可以先分解。例如
\begin{equation*}
 \frac1{(x+1)(x+2)}=\frac1{x+1}-\frac1{x+2},
\end{equation*}
从而原积分化为两个对数。不可约二次因子先配方，则可以与反正切公式联系。这里的代数整理和几何参数选择，分别延续了第四章对二次表达式的研究。

\begin{example}
求\(\int dx/(x^2+2x+5)\)、\(\int dx/(x^2+3x+2)\)及\(\int\sqrt{a^2-x^2}\,dx\)。对最后一式，比较三角代换与圆的几何关系。
\end{example}

\begin{exercise}
\mbox{}\par\nobreak

（1）计算\(\int_0^{a/2} x^2/\sqrt{a^2-x^2}\,dx\)，其中\(a>0\)，并解释所选参数区间。

（2）令\(x=\sinh t\)，求\(\int\sqrt{1+x^2}\,dx\)，并对结果求导核对。

（3）计算\(\int_0^1dx/(x^2+4x+5)\)及\(\int_0^1dx/[(x+1)(x+3)]\)，比较配方与部分分式所利用的代数结构。

（4）对\(a>0\)，用\(x=a\sinh t\)求\(\int dx/\sqrt{x^2+a^2}\)，再将结果写成对数形式，并求导验证。

（5）计算\(\int_0^a(a^2-x^2)^{3/2}\,dx\)，其中\(a>0\)。利用三角代换化为三角函数幂的积分，并与本节分部积分中的递推关系联系起来。
\end{exercise}
